\documentclass[10pt,a4paper]{amsart}
\usepackage[margin=19mm]{geometry}
\usepackage{amsmath,amssymb,mathtools}
\usepackage[scr=rsfs,cal=euler]{mathalfa}
\usepackage{xfrac}
\usepackage[unicode,hidelinks,bookmarks=false]{hyperref}
\newcommand{\ie}{i.e.\ }
\newcommand{\cf}{cf.\ }
\newcommand{\resp}{respectively\ }
\newcommand{\Subsection}{\subsection}
\providecommand{\nobreakdash}{\nobreak\hbox{-}}
\setlength{\emergencystretch}{2em}
\multlinegap0pt
\usepackage{amssymb}
\newtheorem{theorem}{Theorem}
\title{A unification of Euler-Maclaurin and Euler-Boole summation formulas}
\author{Yuan He}
\date{Source: arXiv:2609.00687v1 (CC BY 4.0)}
\begin{document}
\maketitle

\noindent\emph{Source and licence.} A unification of Euler-Maclaurin and Euler-Boole summation formulas. Version arXiv:2609.00687v1 (CC BY 4.0), distributed by arXiv under the Creative Commons Attribution 4.0 International licence, http://creativecommons.org/licenses/by/4.0/, as recorded in the arXiv licence metadata for this exact version and shown on the abstract page https://arxiv.org/abs/2609.00687v1. Full licence notice: \texttt{licenses/arxiv-2609.00687-CC-BY-4.0.txt}.\par\smallskip

\noindent\textit{Editorial scope.} Counted unit: the article's theorem thm2.6 (source0.tex lines 272--284). The article's complete original proof of it is delivered with it (source0.tex lines 286--379, one \texttt{proof} environment of 94 lines). No mathematics is written, repaired or re-derived: the statement and the proof are the article's own lines, in the article's own notation, and no retained line is substituted or re-flowed.  Retained uncounted as the source-local context the proof needs: lines 115--136; lines 132--135; lines 237--242; lines 244--249.  Nothing was omitted from any retained span, so the package carries no omission record.  The retained lines cite 2 works, whose entries are transcribed verbatim from the article's own bibliography source in the e-print and delivered with short editorial labels recorded in the item's reference mapping.  No retained line contains a figure environment, an image-inclusion command or drawing code, so no image is delivered and the package declares no asset dependency.  Editorial formatting only, in addition: an AMS \texttt{amsart} preamble that restates the article's own declarations and macros used by the retained lines, namely the environments \texttt{theorem} on separate counters, together with the standard packages those lines need.  The retained environments print in this document as Theorem 1 in order of appearance, whereas the article prints the same items with the numbering of its own full document.  The complete native source of the article as distributed in the arXiv e-print for this exact version is bundled unchanged as \texttt{sources/209051/source0.tex}.
\begin{theorem}\label{thm2.1} Let $q\in\mathbb{N}$, $m\in\mathbb{N}_{0}$ and $a,b\in\mathbb{R}$ with $a<b$. Let $\chi$ be a Dirichlet character modulo $q$ and let $f(t)$ be an $m+1$
times continuously differentiable function on the interval $[a,b]$. Then, for $\lambda\in\mathbb{C}\setminus\{0\}$,
\begin{eqnarray}\label{eq2.1}
&&\sum_{\substack{a<k\leqslant b\\k\in\mathbb{Z}}}\chi(k)\lambda^{-k}f(k)\nonumber\\
&&\qquad=\sum_{k=0}^{m}\frac{(-1)^{k+1}}{(k+1)!}\bigl(\beta_{k+1,\chi^{-}}(\{b\}_{\chi},\lambda)f^{(k)}(b)-\beta_{k+1,\chi^{-}}(\{a\}_{\chi},\lambda)f^{(k)}(a)\bigl)\nonumber\\
&&\qquad\quad+\delta_{1,\lambda^{q}}\beta_{0,\chi^{-}}(\lambda)\int_{a}^{b}f(t)dt\nonumber\\
&&\qquad\quad+\frac{(-1)^{m}}{(m+1)!}\int_{a}^{b}\beta_{m+1,\chi^{-}}(\{t\}_{\chi},\lambda)f^{(m+1)}(t)dt,
\end{eqnarray}
where
\begin{equation*}
\beta_{0,\chi^{-}}(\lambda)=\frac{1}{q}\sum_{r=1}^{q}\chi(-r)\lambda^{r},
\end{equation*}
and
\begin{equation*}
\beta_{m,\chi^{-}}(\{x\}_{\chi},\lambda)=q^{m-1}\sum_{r=1}^{q}\chi(-r)\lambda^{r-q[\frac{x+r}{q}]}\beta_{m}\biggl(\biggl\{\frac{x+r}{q}\biggl\},\lambda^{q}\biggl),
\end{equation*}
in which $\beta_{n}(x,\lambda)$ denotes the Apostol-Bernoulli polynomials defined for $\lambda\in\mathbb{C}\setminus\{0\}$ by the generating function (see, e.g., \cite{apostol2})
\begin{equation}\label{eq2.2}
\frac{te^{xt}}{\lambda e^{t}-1}=\sum_{n=0}^{\infty}
\beta_{n}(x,\lambda)\frac{t^{n}}{n!}\quad(|t+\log\lambda|<2\pi).
\end{equation}
\end{theorem}

\begin{equation}
\frac{te^{xt}}{\lambda e^{t}-1}=\sum_{n=0}^{\infty}
\beta_{n}(x,\lambda)\frac{t^{n}}{n!}\quad(|t+\log\lambda|<2\pi).
\end{equation}

\begin{equation}\label{eq2.10}
\beta_{0}(x,\lambda)=\begin{cases}
1,  &\lambda=1,\\
0,  &\lambda\not=1,
\end{cases}
\end{equation}

\begin{equation}\label{eq2.11}
\beta_{1}(x,\lambda)=\begin{cases}
x-\frac{1}{2},  &\lambda=1,\\
\frac{1}{\lambda-1},  &\lambda\not=1.
\end{cases}
\end{equation}

\begin{theorem}\label{thm2.6} Let $n,q\in\mathbb{N}$ and $x,a\in\mathbb{R}$ with $a>0$. Let $\chi$ be a Dirichlet character modulo $q$. If $q=1$, then
\begin{equation}\label{eq2.15}
L(1-n,x,a,\chi)=-\frac{\beta_{n}(a,e^{\frac{2\pi \mathrm{i}x}{q}})}{n},
\end{equation}
and if $q\geqslant2$, then
\begin{equation}\label{eq2.16}
L(1-n,x,a,\chi)=-\frac{\beta_{n,\chi}(a,e^{\frac{2\pi \mathrm{i}x}{q}})}{n},
\end{equation}
where $\beta_{n,\chi}(x,\lambda)$ are the generalized Apostol-Bernoulli polynomials, defined for $\lambda\in\mathbb{C}\setminus\{0\}$ and for a Dirichlet character $\chi$ modulo $q\in\mathbb{N}$, by the generating function (see, e.g., \cite[Definition 2.1]{he})
\begin{equation}\label{eq2.17}
\sum_{r=1}^{q}\frac{\chi(r)\lambda^{r}te^{(r+x)t}}{\lambda^{q}e^{qt}-1}=\sum_{n=0}^{\infty}\beta_{n,\chi}(x,\lambda)\frac{t^{n}}{n!}\quad(|t+\log\lambda|<2\pi/q).
\end{equation}
\end{theorem}

\begin{proof} Taking $a=0$, $b=n$ and $f(t)=(t+y)^{s}$, and replacing $m$ by $m-1$ in Theorem \ref{thm2.1}, we find that for $m\in\mathbb{N}$, $n\in\mathbb{N}_{0}$, $y\in\mathbb{R}$ with $y>0$, $\lambda\in\mathbb{C}\setminus\{0\}$, and $s\in\mathbb{C}\setminus\{-1\}$,
\begin{eqnarray}\label{eq2.18}
&&\sum_{k=0}^{n}\chi(k)\lambda^{-k}(k+y)^{s}\nonumber\\
&&\qquad=\sum_{k=2}^{m}\binom{s}{k-1}\frac{(-1)^{k}}{k}\beta_{k,\chi^{-}}(\{n\}_{\chi},\lambda)(n+y)^{s+1-k}\nonumber\\
&&\qquad\quad-\sum_{k=2}^{m}\binom{s}{k-1}\frac{(-1)^{k}}{k}\beta_{k,\chi^{-}}(\{0\}_{\chi},\lambda)y^{s+1-k}\nonumber\\
&&\qquad\quad+\chi(0)y^{s}-\beta_{1,\chi^{-}}(\{n\}_{\chi},\lambda)(n+y)^{s}+\beta_{1,\chi^{-}}(\{0\}_{\chi},\lambda)y^{s}\nonumber\\
&&\qquad\quad+\delta_{1,\lambda^{q}}\beta_{0,\chi^{-}}(\lambda)\biggl(\frac{(n+y)^{s+1}}{s+1}-\frac{y^{s+1}}{s+1}\biggl)\nonumber\\
&&\qquad\quad+(-1)^{m-1}\binom{s}{m}\int_{0}^{n}\beta_{m,\chi^{-}}(\{t\}_{\chi},\lambda)(t+y)^{s-m}dt,
\end{eqnarray}
where the summation on the right hand side of \eqref{eq2.18} vanishes when $m=1$, and $\binom{s}{k}$ are the generalized binomial coefficients defined by
\begin{equation*}
\binom{s}{0}=1,\quad\binom{s}{k}=\frac{s(s-1)\cdots(s-k+1)}{k!}\quad(k\in\mathbb{N}).
\end{equation*}
Let $\overline{\beta}_{n,\chi}(x,\lambda)$ be the $n$-th generalized Apostol-Bernoulli function defined (see, e.g., \cite[Definition 2.7]{he}):
\begin{equation}\label{eq2.19}
\overline{\beta}_{n,\chi}(x,\lambda)=\beta_{n,\chi}(\{x\}_{\chi},\lambda)+\frac{1}{2}\delta_{1,n}\delta_{\mathbb{Z}}(x)\lambda^{-x}\chi(-x),
\end{equation}
where $\beta_{n,\chi}(\{x\}_{\chi},\lambda)=\chi(-1)\beta_{n,\chi^{-}}(\{x\}_{\chi},\lambda)$, $\delta_{\mathbb{Z}}(x)=1$ if $x\in\mathbb{Z}$ and $0$ otherwise. This function is shown in \cite[Theorem 3.1]{he} to have the Fourier series
\begin{equation}\label{eq2.20}
\overline{\beta}_{n,\chi}(x,\lambda)=-\frac{q^{n-1}n!}{\lambda^{x}(2\pi \mathrm{i})^{n}}\sideset{}{'}\sum_{k=-\infty}^{+\infty}\frac{G(k,\chi)e^{\frac{2\pi\mathrm{i}kx}{q}}}{(k-\frac{q\log\lambda}{2\pi\mathrm{i}})^{n}},
\end{equation}
where the dash denotes throughout that undefined terms are excluded from the sum, and $G(k,\chi)$ is the Gauss sum defined by
\begin{equation*}
G(k,\chi)=\sum_{r=1}^{q}\chi(r)e^{\frac{2\pi \mathrm{i}rk}{q}}.
\end{equation*}
Hence, we see from \eqref{eq2.11}, \eqref{eq2.19} and \eqref{eq2.20} that if $|\lambda|=1$ and $\Re(s)<m-1$, then the integral
\begin{equation*}
R(n)=\int_{n}^{\infty}\beta_{m,\chi^{-}}(\{t\}_{\chi},\lambda)(t+y)^{s-m}dt
\end{equation*}
converges absolutely, and $|R(n)|\rightarrow0$ as $n\rightarrow\infty$. So, when $|\lambda|=1$, $\Re(s)<m-1$ and $s\neq-1$, we can write \eqref{eq2.18} as
\begin{eqnarray}\label{eq2.21}
&&\sum_{k=0}^{n}\chi(k)\lambda^{-k}(k+y)^{s}\nonumber\\
&&\qquad=\sum_{k=2}^{m}\binom{s}{k-1}\frac{(-1)^{k}}{k}\beta_{k,\chi^{-}}(\{n\}_{\chi},\lambda)(n+y)^{s+1-k}\nonumber\\
&&\qquad\quad-\sum_{k=2}^{m}\binom{s}{k-1}\frac{(-1)^{k}}{k}\beta_{k,\chi^{-}}(\{0\}_{\chi},\lambda)y^{s+1-k}\nonumber\\
&&\qquad\quad+\chi(0)y^{s}-\beta_{1,\chi^{-}}(\{n\}_{\chi},\lambda)(n+y)^{s}+\beta_{1,\chi^{-}}(\{0\}_{\chi},\lambda)y^{s}\nonumber\\
&&\qquad\quad+\delta_{1,\lambda^{q}}\beta_{0,\chi^{-}}(\lambda)\biggl(\frac{(n+y)^{s+1}}{s+1}-\frac{y^{s+1}}{s+1}\biggl)\nonumber\\
&&\qquad\quad+(-1)^{m-1}\binom{s}{m}\int_{0}^{\infty}\beta_{m,\chi^{-}}(\{t\}_{\chi},\lambda)(t+y)^{s-m}dt\nonumber\\
&&\qquad\quad+(-1)^{m}\binom{s}{m}\int_{n}^{\infty}\beta_{m,\chi^{-}}(\{t\}_{\chi},\lambda)(t+y)^{s-m}dt.
\end{eqnarray}
In particular, for $\Re(s)<-1$, letting $n\rightarrow\infty$ in \eqref{eq2.21}, we obtain
\begin{eqnarray}\label{eq2.22}
&&\sum_{k=0}^{\infty}\frac{\chi(k)\lambda^{-k}}{(k+y)^{-s}}\nonumber\\
&&\qquad=-\sum_{k=2}^{m}\binom{s}{k-1}\frac{(-1)^{k}}{k}\beta_{k,\chi^{-}}(\{0\}_{\chi},\lambda)y^{s+1-k}\nonumber\\
&&\qquad\quad+\chi(0)y^{s}+\beta_{1,\chi^{-}}(\{0\}_{\chi},\lambda)y^{s}-\delta_{1,\lambda^{q}}\beta_{0,\chi^{-}}(\lambda)\frac{y^{s+1}}{s+1}\nonumber\\
&&\qquad\quad+(-1)^{m-1}\binom{s}{m}\int_{0}^{\infty}\beta_{m,\chi^{-}}(\{t\}_{\chi},\lambda)(t+y)^{s-m}dt.
\end{eqnarray}
Inserting \eqref{eq2.22} into \eqref{eq2.21}, we know that for $|\lambda|=1$,
\begin{eqnarray}\label{eq2.23}
&&\sum_{k=0}^{n}\chi(k)\lambda^{-k}(k+y)^{s}\nonumber\\
&&\qquad=\sum_{k=0}^{\infty}\frac{\chi(k)\lambda^{-k}}{(k+y)^{-s}}+\sum_{k=2}^{m}\binom{s}{k-1}\frac{(-1)^{k}}{k}\beta_{k,\chi^{-}}(\{n\}_{\chi},\lambda)(n+y)^{s+1-k}\nonumber\\
&&\qquad\quad-\beta_{1,\chi^{-}}(\{n\}_{\chi},\lambda)(n+y)^{s}+\delta_{1,\lambda^{q}}\beta_{0,\chi^{-}}(\lambda)\frac{(n+y)^{s+1}}{s+1}\nonumber\\
&&\qquad\quad+(-1)^{m}\binom{s}{m}\int_{n}^{\infty}\beta_{m,\chi^{-}}(\{t\}_{\chi},\lambda)(t+y)^{s-m}dt,
\end{eqnarray}
which holds for $\Re(s)\leqslant m-1$ and $s\neq-1$ by analytic continuation. Now, taking $s=m-1$ and $n=0$ in \eqref{eq2.23}, we have
\begin{eqnarray}\label{eq2.24}
&&\sum_{k=0}^{\infty}\frac{\chi(k)\lambda^{-k}}{(k+y)^{1-m}}\nonumber\\
&&\qquad=-\frac{1}{m}\sum_{k=0}^{m}\binom{m}{k}(-1)^{k}\beta_{k,\chi^{-}}(\{0\}_{\chi},\lambda)y^{m-k}+\chi(0)y^{m-1}.
\end{eqnarray}
Note that from \eqref{eq2.2} and \eqref{eq2.17} we have
\begin{equation}\label{eq2.25}
\beta_{n,\chi}(x,\lambda)=q^{n-1}\sum_{r=1}^{q}\chi(r)\lambda^{r}\beta_{n}\biggl(\frac{x+r}{q},\lambda^{q}\biggl)\quad(n\in\mathbb{N}_{0}).
\end{equation}
Moreover, for $n\in\mathbb{N}$ and $x\in\mathbb{R}$ (see, e.g., \cite[Propostion 2.8]{he}),
\begin{equation}\label{eq2.26}
\beta_{n,\chi}(\{x\}_{\chi},\lambda)=\beta_{n,\chi}(x,\lambda)-n\underset{\substack{0\leqslant r\leqslant x\\ r\in\mathbb{Z}}}{\sum}\chi(-r)\lambda^{-r}(x-r)^{n-1}.
\end{equation}
From \eqref{eq2.10}, \eqref{eq2.25} and \eqref{eq2.26}, it follows that for $n\in\mathbb{N}_{0}$,
\begin{equation}\label{eq2.27}
\beta_{n,\chi^{-}}(\{0\}_{\chi},\lambda)=\beta_{n,\chi^{-}}(0,\lambda)-\delta_{1,n}\chi(0).
\end{equation}
Applying \eqref{eq2.27} to the right-hand side of \eqref{eq2.24}, we have
\begin{equation}\label{eq2.28}
\sum_{k=0}^{\infty}\frac{\chi(k)\lambda^{-k}}{(k+y)^{1-m}}=-\frac{1}{m}\sum_{k=0}^{m}\binom{m}{k}(-1)^{k}\beta_{k,\chi^{-}}(0,\lambda)y^{m-k}.
\end{equation}
It is shown in \cite[Proposition 2.3]{he} that if $q=1$, then
\begin{equation}\label{eq2.29}
\beta_{n,\chi}(-x,\lambda)=(-1)^{n}\beta_{n}(x,\lambda^{-1})\quad(n\in\mathbb{N}_{0}),
\end{equation}
and if $q\geqslant2$, then
\begin{equation}\label{eq2.30}
\beta_{n,\chi}(-x,\lambda)=(-1)^{n}\beta_{n,\chi^{-}}(x,\lambda^{-1})\quad(n\in\mathbb{N}_{0}).
\end{equation}
Thus, replacing $\lambda$ by $\lambda^{-1}$ in \eqref{eq2.28}, with the help of \eqref{eq2.29} and \eqref{eq2.30}, we find from \eqref{eq2.2} and \eqref{eq2.17} that if $q=1$, then
\begin{eqnarray*}
\sum_{k=0}^{\infty}\frac{\chi(k)\lambda^{k}}{(k+y)^{1-m}}&=&-\frac{1}{m}\sum_{k=0}^{m}\binom{m}{k}\beta_{k}(0,\lambda)y^{m-k}\nonumber\\
&=&-\frac{\beta_{m}(y,\lambda)}{m},
\end{eqnarray*}
and if $q\geqslant2$, then
\begin{eqnarray*}
\sum_{k=0}^{\infty}\frac{\chi(k)\lambda^{k}}{(k+y)^{1-m}}&=&-\frac{1}{m}\sum_{k=0}^{m}\binom{m}{k}\beta_{k,\chi}(0,\lambda)y^{m-k}\nonumber\\
&=&-\frac{\beta_{m,\chi}(y,\lambda)}{m}.
\end{eqnarray*}
Now the desired result follows by taking $\lambda=e^{\frac{2\pi \mathrm{i}x}{q}}$ and replacing $y$ by $a$ and $m$ by $n$ in the above two identities. This completes the proof of Theorem \ref{thm2.6}.
\end{proof}

\begin{thebibliography}{99}
\bibitem[aposto]{apostol2} T. M. Apostol, On the Lerch zeta function, \emph{Pacific J. Math.}, \textbf{1} (1951), 161--167.
\bibitem[he]{he} Y. He, On the special values of Berndt's generalized $L$-function, arXiv:2608.29636.
\end{thebibliography}
\end{document}
